% Pursuing Stacks (cotes 134-2 to 134-8): archive cote 134-7, site page 43 (PDF page 44), typescript page 485.
% typescript. Draft compared: ps6.tex:2108-2184.
% AI-transcribed by Claude Opus 5.5 (claude-opus-5-5, adaptive thinking, effort high), Message Batches API, batch ps3 (msgbatch_01XdFmuVVBAcgc7EFxdtCqsj), 2026-11-12; not yet checked by eye.
\scanpage{134-7/43}
\mspage{485}
functor (7'), this is clearly false, as we do not have any canonical
map from W(M) into $W(\Gamma_k^{\wedge}(M))$ , only into $W^{\wedge}\Gamma_k^{\wedge}(M)$ \xout{wh} - we have an
embedding
\[ (12) \qquad W(\Gamma_k^{\wedge}(M)) \hand{\hookrightarrow} W^{\wedge}\Gamma_k^{\wedge}(M) \quad , \]
but it is clear that in general, the exponential map (11) does not
factor through the first term in (12). (It does of course when we look
at sections over k only, but when we go over to a general k', we hit
into the trouble that formation of inverse limits des not commute with
ring extension $\otimes_k k'$ !). We may however express (9) by stating that the
functor (7') is ``\uline{prorepresentable}'' by the \uline{pro}-\uline{object}
\[ (13) \qquad \mathrm{Pr}\bullet\, \Gamma_k(M) \overset{\mathrm{def}}{=} (\Gamma_k(M)(i))_{i \geqslant 0} \quad \text{in} \quad \mathrm{Pro}(\mathrm{Ab}_k) \quad , \]
this is even a \uline{strict} proobject (the transition morphisms are epimor-%
phisms), which implies that the functor it prorepresents is represen-%
table iff this projective system is ``essentially constant'' in the most
trivial sense, which means here
\[ \Gamma_k^{\,i}(M) \;=\; 0 \quad \text{for large } i \quad , \]
a condition which presumably is satisfied only for $M=0$ ! Thus, the
``correct'' interpretation of non-pointed \inserted{(quasi-coherent)} linearization seems to me to
be the corresponding functor, which I would like now to call $L_k$ or
simply L as before \hand{but with slightly different meaning:}
\[ (14) \qquad L \text{ or } L_k \;:\; U(k) \hand{\longrightarrow} \mathrm{Pr}\bullet(\mathrm{Ab}_k) \hand{\underline{\approx}} \mathrm{Pr}\bullet(\mathrm{Qucoh}(k)) \quad , \]
where \xout{Qu} Qucoh(k) is defined in (8). In computational terms, I
would like to view \xout{Qucohx}\inserted{L(X)} (for a unipotent bundle X) \deleted{\hand{\unsure{L(X)}}} to be a
pro-k-module, but in terms of geometric intuition, I would see it
rather as a pro-$\underline{O}_k$-Module, i.e. essentially as an inverse system of
quasicoherent modules. It is in these latter terms that the constructi-%
on we just gave generalises to unipotent bundles over arbitrary ground
schemes, not necessarily affine ones. As for the ``pointed'' quasicohe-%
rent linearization functor
\[ (15) \qquad L_{pt} \text{ or } L_{k\ pt} \;:\; U(k)^{\bullet} \hand{\longrightarrow} (\mathrm{Ab}_k) \hand{\underline{\approx}} \mathrm{Qucoh}(k) \quad , \]
which I like best to view as taking values $L_{pt}(X)$ which are quasicoherent
sheaves, it maps into L by
\[ (1\over{\unsure{5}}{6}) \qquad L_{pt}(X) \hand{\hookrightarrow} L(X) \quad , \]
interpreting objects of \xout{categoriexx} a category as special cases of pro-%
objects. We'll denote by
\[ \hand{(17)} \qquad WL(X) \hand{\in} \mathrm{Pro}(U(k)) \]
\margin{43}
